\documentclass[10pt,a4paper]{amsart}
\usepackage[margin=19mm]{geometry}
\usepackage{amsmath,amssymb,mathtools}
\usepackage[scr=rsfs,cal=euler]{mathalfa}
\usepackage{xfrac}
\usepackage[unicode,hidelinks,bookmarks=false]{hyperref}
\newcommand{\ie}{i.e.\ }
\newcommand{\cf}{cf.\ }
\newcommand{\resp}{respectively\ }
\newcommand{\Subsection}{\subsection}
\providecommand{\nobreakdash}{\nobreak\hbox{-}}
\setlength{\emergencystretch}{2em}
\multlinegap0pt
\usepackage{bm}
\usepackage{bbm}
\usepackage{dsfont}
\usepackage{xspace}
\usepackage{cancel}
\usepackage{mathtools}
\usepackage{amsthm}
\makeatletter
\@ifundefined{theorem}{\newtheorem{theorem}{Theorem}}{}
\@ifundefined{definition}{\newtheorem{definition}[theorem]{Definition}}{}
\@ifundefined{lemma}{\newtheorem{lemma}[theorem]{Lemma}}{}
\makeatother
\newcommand{\AGL}{\mathrm{AGL}_1}
\newcommand{\comm}[2]{[#1,#2]}
\title[Cross-Commuting Nonabelian Squares in Affine Groups over Fin]{Cross-Commuting Nonabelian Squares in Affine Groups over Finite Commutative Principal Ideal Rings}
\author{Kenta Kasai}
\date{Source: arXiv:2604.02063v1 (CC BY 4.0)}
\begin{document}
\maketitle

\noindent\emph{Source and licence.} Cross-Commuting Nonabelian Squares in Affine Groups over Finite Commutative Principal Ideal Rings. Version arXiv:2604.02063v1 (CC BY 4.0), distributed under the Creative Commons Attribution 4.0 International licence, as recorded in the arXiv licence metadata for this exact version and shown on the abstract page https://arxiv.org/abs/2604.02063v1. Full rights notice: licenses/arxiv-2604.02063-CC-BY-4.0.txt.\par\smallskip

\noindent\textit{Editorial scope.} Counted unit: the Theorem whose source label is \texttt{thm:local-obstruction-en} in the article (source0.tex lines 171--180, source label \texttt{thm:local-obstruction-en}), delivered together with the article's own complete original proof of it (source0.tex lines 182--250, one \texttt{proof} environment of 69 lines). No mathematics is written, repaired or re-derived: statement and proof are the article's own text line for line. Required context retained uncounted: lem:comm-criterion-en (lines 122--128); def:square-en (lines 141--152). Every definition, display and label of those lines is retained unchanged. Omitted from this package and not redistributed: 1--121, 129--140, 153--170, 181--181, 251--606. Nothing inside a retained excerpt is omitted or altered: every retained body is delivered exactly as the article's lines are written, so no omission receipt of any kind accompanies this package. Citations: the retained excerpts carry no citation command at all, so no bibliography entry of the article is transcribed and no reference is redistributed. Reference adaptation: none was needed and none was made; every reference inside the delivered unit resolves inside it, so no unresolved reference or citation is produced. Printed slips kept literal: no printed word, symbol, hypothesis or number of the counted statement or of its proof was altered. Layout and class: the article's own class is \texttt{article} with packages including geometry, amsmath, amssymb, amsthm, mathtools, bm, hyperref; the wrapper is the lane's pinned \texttt{amsart} editorial wrapper at 10pt on A4, which restates the theorem-like environments the retained text needs and adds the article's own macro definitions that the retained text needs (\texttt{\textbackslash AGL}, \texttt{\textbackslash comm}), with extra wrapper packages (bm, bbm, dsfont, xspace, cancel, mathtools). The delivered numbering is the unit's own. Assets: no retained span contains a figure environment, a graphics-inclusion command, a PDF-page inclusion, a table or a TikZ picture; no image file is copied into this package, the package declares no asset dependency and no image file is redistributed with it. Licence: version arXiv:2604.02063v1 of this article is distributed under the Creative Commons Attribution 4.0 International licence, as recorded in the arXiv licence metadata for this exact version and shown on the abstract page https://arxiv.org/abs/2604.02063v1. A full rights notice is delivered at licenses/arxiv-2604.02063-CC-BY-4.0.txt. Characters: every retained excerpt is pure ASCII, so the delivered engine, XeLaTeX with its default Latin Modern text font, reports no missing character.
\begin{lemma}[Commutation criterion]
\label{lem:comm-criterion-en}
Two affine permutations $(a,b)$ and $(c,d)$ in $\AGL(R)$ commute if and only if
\[
(a-1)d=(c-1)b.
\]
\end{lemma}

\begin{definition}[Cross-commuting nonabelian square]
\label{def:square-en}
A quadruple $(F_0,F_1,G_0,G_1)$ in $\AGL(R)$ is called a
\emph{cross-commuting nonabelian square} if
\[
\comm{F_i}{G_j}=1\qquad (i,j\in\{0,1\}),
\]
but
\[
\comm{F_0}{F_1}\ne 1,\qquad \comm{G_0}{G_1}\ne 1.
\]
\end{definition}

\begin{theorem}[Local obstruction]
\label{thm:local-obstruction-en}
Let $R$ be a finite commutative local principal ideal ring. If
$F_0,F_1\in\AGL(R)$ do not commute, then the common centralizer
\[
C(F_0)\cap C(F_1)
\]
is abelian. In particular, no cross-commuting nonabelian square exists in
$\AGL(R)$.
\end{theorem}

\begin{proof}
Write
\[
F_0=(a_0,b_0),\qquad F_1=(a_1,b_1),
\]
and set
\[
s_0:=a_0-1,\qquad s_1:=a_1-1.
\]
By Lemma~\ref{lem:comm-criterion-en}, the noncommutativity of $F_0$ and $F_1$
means that
\[
\Delta:=s_0b_1-s_1b_0
\]
is nonzero.

Now let $G=(c,d)$ commute with both $F_0$ and $F_1$, and write $u:=c-1$. By
Lemma~\ref{lem:comm-criterion-en}, the pair $(u,d)$ satisfies
\[
b_0u-s_0d=0,\qquad b_1u-s_1d=0.
\]
Thus $(u,d)^{\mathsf T}$ lies in the kernel of the matrix
\[
M:=
\begin{pmatrix}
b_0 & -s_0\\
b_1 & -s_1
\end{pmatrix}.
\]
Its determinant is exactly $\Delta$.

Since $R$ is a principal ideal ring, $M$ admits a Smith normal form:
there exist invertible matrices $U,V\in \operatorname{GL}_2(R)$ and integers
$0\le \alpha\le \beta\le \ell$ such that
\[
UMV=\operatorname{diag}(\pi^\alpha,\pi^\beta),
\qquad
\alpha+\beta=v_\pi(\Delta)<\ell.
\]
Hence
\[
\ker(M)=V\bigl(\pi^{\ell-\alpha}R\times \pi^{\ell-\beta}R\bigr).
\]

Take two elements $G=(1+u,d)$ and $G'=(1+u',d')$ in $C(F_0)\cap C(F_1)$. Their
parameter vectors lie in $\ker(M)$, so there exist $r,s,r',s'\in R$ with
\[
\binom{u}{d}
=
V\binom{\pi^{\ell-\alpha}r}{\pi^{\ell-\beta}s},
\qquad
\binom{u'}{d'}
=
V\binom{\pi^{\ell-\alpha}r'}{\pi^{\ell-\beta}s'}.
\]
Therefore
\[
ud'-u'd
=
\det(V)\,\pi^{2\ell-\alpha-\beta}(rs'-r's).
\]
Because $\alpha+\beta<\ell$, we have $2\ell-\alpha-\beta>\ell$, and therefore
$\pi^{2\ell-\alpha-\beta}=0$. Hence $ud'=u'd$. By
Lemma~\ref{lem:comm-criterion-en}, $G$ and $G'$ commute.

Thus every two elements of $C(F_0)\cap C(F_1)$ commute, so this common
centralizer is abelian. The final assertion follows immediately from
Definition~\ref{def:square-en}.
\end{proof}

\end{document}
